
\clearpage
\section{FUNDAMENTAÇÃO LEGAL}
\label{FUNDAMENTAÇÃO LEGAL}

No processo de elaboração e atualização deste projeto pedagógico, foram consideradas as normativas legais em âmbito nacional e institucional que regulamentam os cursos superiores de graduação, especialmente as normativas específicas para os cursos de bacharelado, a saber:

\subsection{Normativas Nacionais}

\begin{itemize}
	\item Lei Nº 9.394, de 20 de dezembro de 1996, que estabelece as Diretrizes e Bases da Educação Nacional.
	\item Lei Nº 11.892, de 29 de dezembro de 2008, que institui a Rede Federal de Educação Profissional, Científica e Tecnológica, cria os Institutos Federais e dá outras providências.
	\item Parecer CES Nº 277/2006, que versa sobre nova forma de organização da Educação Profissional e Tecnológica de graduação.
	\item Parecer CNE/CES Nº 583, de 4 de abril de 2001, que dispõe sobre a orientação para as Diretrizes Curriculares dos Cursos de Graduação.
	\item Parecer CNE/CES nº 1.362/2001, aprovado em 12 de dezembro de 2001 que trata sobre as Diretrizes Curriculares Nacionais dos Cursos de Engenharia.
	\item Resolução CNE/CES nº 11, de 11 de março de 2002 que institui as Diretrizes Curriculares Nacionais do Curso de Graduação em Engenharia.
	\item Parecer CNE/CES Nº 8/2007, de 31 de janeiro de 2007, que dispõe sobre carga horária mínima e procedimentos relativos à integralização e duração dos cursos de graduação, bacharelados, na modalidade presencial. 
	\item Resolução CNE/CES Nº 2, de 18 de junho de 2007, que dispõe sobre carga horária mínima e procedimentos relativos à integralização e duração dos cursos de graduação, bacharelados, na modalidade presencial. 
	\item Resolução CNE/CES Nº 3, de 2 de julho de 2007, dispõe sobre procedimentos a serem adotados quanto ao conceito de hora-aula, e dá outras providências. 
	\item Lei Nº 10.861, de 14 de abril de 2004, que institui o Sistema Nacional de Avaliação da Educação Superior – SINAES e dá outras providências. 
	\item Decreto Nº 5.773, de 9 de maio de 2006, que dispõe sobre o exercício das funções de regulação, supervisão e avaliação de instituições de educação superior e cursos superiores de graduação e sequenciais no sistema federal de ensino. 
	\item Portaria MEC Nº 40, de 12 de dezembro de 2007, reeditada em 29 de dezembro de 2011, que institui o e -MEC, sistema eletrônico de fluxo de trabalho e gerenciamento de informa ações relativas aos processos de regulação, avaliação e supervisão da educação superior no sistema federal de educação, e o Cadastro e - MEC de Instituições e Cursos Superiores e consolida disposições sobre indicadores de qualidade, banco de avaliadores (Basis) e o Exame Nacional de Desempenho de Estudantes (Enade) e outras disposições. 
	\item Lei nº 9.795, de 27 de abril de 1999 e Decreto Nº 4.281 de 25 de junho de 2002, que tratam sobre as Políticas de educação ambiental.  
	\item Resolução CNE/CP Nº 2, de 15 de junho de 2012, que estabelece as Diretrizes Curriculares Nacionais para a Educação Ambiental. 
	\item Parecer CNE/CP Nº 8, de 06 de março de 2012 e Resolução CNE/CP Nº 1, de 30 de maio de 2012, que tratam sobre as Diretrizes Nacionais para a Educação em Direitos Humanos. 
	\item Resolução CNE/CP Nº 1, de 17 de junho de 2004, que institui Diretrizes Curriculares Nacionais para a Educação das Relações Étnico-Raciais e para o Ensino de História e Cultura Afro-Brasileira e Africana. 
	\item Decreto Nº 5.626, de 22 de dezembro de 2005, que regulamenta a Lei Nº 10.436, de 24 de abril de 2002, que dispões sobre a Língua Brasileira de Sinais – Libras. 
	\item Lei N° 12.764, de 27 de dezembro de 2012, que trata da Proteção dos Direitos da Pessoa com Transtorno do Espectro Autista. 
	\item Condições de acessibilidade para pessoas com deficiência ou mobilidade reduzida, conforme disposto na CF/88, art. 205, 206 e 208, na NBR 9050/2004, da ABNT, na Lei N° 10.098/2000, nos Decretos N° 5.296/2004, N° 6.949/2009, N° 7.611/2011 e na Portaria N° 3.284/2003.
	\item Resolução CONAES N°12/2016, que normatiza o Núcleo Docente Estruturante e dá outras providências;
	\item Instrumentos para autorização, renovação e reconhecimento dos cursos, publicados pelo INEP.
\end{itemize}

\subsection{Normativas Institucionais}

\begin{itemize}
	\item Regulamento da Organização Didática no IFCE – ROD. 
	\item Plano de Desenvolvimento Institucional do IFCE – PDI. 
	\item Projeto Pedagógico Institucional – PPI. 
	\item Tabela de Perfil Docente. 
	\item Resolução CONSUP Nº 028, de 08 de agosto de 2014, que dispõe sobre o Manual de Estágio do IFCE. 
	\item Resolução que regulamenta a Carga Horária docente. 
	\item Resolução CONSUP Nº 004, de 28 de janeiro de 2015, que determina a organização do Núcleo Docente Estruturante no IFCE. 
	\item Resolução CONSUP Nº 099, de 27 de setembro de 2017, que aprova o Manual de Elaboração de Projetos Pedagógicos de Cursos do IFCE.
	\item Resolução CONSUP Nº 75, de 13 de agosto de 2018, que define as normas de funcionamento do colegiado dos cursos técnicos e de graduação do IFCE.
\end{itemize}
